\chapter{Literature Review Protocol}
\label{app:master-review-protocol}

This appendix records version 1.0 of the structured literature mapping
protocol, frozen on 21 June 2026. It is the operational companion to the
methodology summarized in \cref{ch:master-background}. The review follows the
planning, selection, quality-assessment, extraction, and synthesis principles
of software-engineering review guidance \cite{kitchenham2007guidelines}, while
explicitly acknowledging that the seed corpus was assembled iteratively before
the protocol was frozen.

\section{Research Questions}

\begin{description}
  \item[RQ-R1] Which evidence sources are used to identify Android
    third-party components from compiled applications?
  \item[RQ-R2] Which techniques support library identity, exact-version
    attribution, or both?
  \item[RQ-R3] Which compilation, optimization, obfuscation, packaging, and
    dependency transformations threaten each evidence source?
  \item[RQ-R4] How do studies construct ground truth and evaluate
    library-level, version-level, and vulnerability-related claims?
  \item[RQ-R5] How do studies represent confidence, ambiguity, provenance,
    abstention, partial inclusion, and conflicting evidence?
  \item[RQ-R6] Which research gaps prevent recovered component claims from
    becoming trustworthy and auditable \SBOM statements?
\end{description}

\section{Search Strategy}

\subsection{Sources and Scope}

The primary search sources are ACM Digital Library, IEEE Xplore, SpringerLink,
ScienceDirect, USENIX, and arXiv. Google Scholar is used for discovery and
backward or forward citation chaining rather than as the sole bibliographic
authority. The publication window begins in 2015 and ends on the date recorded
for each search refresh. English full text is preferred; English title and
abstract are mandatory for screening.

\subsection{Concept Blocks}

\begin{table}[H]
  \centering
  \caption{Search concept blocks}
  \label{tab:review-search-blocks}
  \begin{tabularx}{\textwidth}{p{0.18\textwidth}X}
    \toprule
    Block & Terms \\
    \midrule
    Platform & Android, APK, DEX, Dalvik \\
    Component & third-party library, TPL, software component, dependency,
      SBOM, software bill of materials \\
    Task & detect*, identif*, recogn*, recover*, attribut*, fingerprint*,
      match* \\
    Threat or specificity & obfuscat*, R8, ProGuard, optim*, shrink*, inlin*,
      version*, vulnerab* \\
    \bottomrule
  \end{tabularx}
\end{table}

The canonical query is:

\begin{quote}\footnotesize\ttfamily
(Android OR APK OR DEX OR Dalvik)\\
AND ("third-party library" OR TPL OR "software component" OR dependency\\
\hspace*{1.5em}OR SBOM OR "software bill of materials")\\
AND (detect* OR identif* OR recogn* OR recover* OR attribut*\\
\hspace*{1.5em}OR fingerprint* OR match*)\\
AND (obfuscat* OR R8 OR ProGuard OR optim* OR shrink* OR inlin*\\
\hspace*{1.5em}OR version* OR vulnerab*)
\end{quote}

Database-specific syntax may change, but a refresh must preserve the four
concept blocks and record the exact executed query, date, raw result count,
deduplicated count, and selection decisions.

The frozen seed corpus predates this audit requirement. Its original
database-specific query executions, raw result counts, and deduplicated flow
are therefore not claimed as retrospectively available.

\section{Inclusion and Exclusion Criteria}

\subsection{Inclusion}

A publication is included when it:

\begin{enumerate}
  \item addresses component, library, package, or version recovery from Android
    binaries or evaluates tools that do;
  \item studies evidence obtainable from compiled artifacts, reference
    libraries, package metadata, or related dependency knowledge;
  \item contributes a method, benchmark, empirical assessment, taxonomy, or
    SBOM-oriented framework relevant to compiled-artifact attribution;
  \item permits extraction of at least the task, evidence source, evaluation
    basis, and limitations; and
  \item has an accessible English title and abstract.
\end{enumerate}

\subsection{Exclusion}

A publication is excluded when it:

\begin{enumerate}
  \item relies only on source dependency manifests and does not study binary
    recovery;
  \item concerns malware classification, clone detection, or vulnerability
    scanning without a component-attribution contribution;
  \item is a tutorial, slide deck, blog post, product page, or duplicate;
  \item reports no separable Android result;
  \item lacks enough information to identify its method or evaluation basis;
    or
  \item is a shorter version of an included extended study and contributes no
    distinct method, dataset, or result.
\end{enumerate}

\subsection{Screening and Evidence Tiers}

Records are deduplicated and screened by title and abstract. Full-text
eligibility screening and structured extraction are completed for Tier A.
Records without a verified local full text remain Tier B pending that stage.
The evidence tiers are:

\begin{description}
  \item[Tier A --- full-text verified] A local full text has been inspected
    and the publication may support structured and quantitative extraction.
  \item[Tier B --- supplementary record] Metadata or a structured note exists,
    but the local full text has not yet been verified. The record may guide
    search and provisional gap mapping, but not precise quantitative claims.
  \item[Tier C --- contextual authority] Standards, official documentation,
    and project artifacts used for definitions or framework design. They are
    not counted as primary studies.
  \item[Excluded] The reason for exclusion is recorded and the item is not
    used in synthesis.
\end{description}

\section{Data Extraction Schema}

\begin{longtable}{p{0.24\textwidth}p{0.68\textwidth}}
  \caption{Structured extraction fields}
  \label{tab:review-extraction-schema}\\
  \toprule
  Dimension & Extracted information \\
  \midrule
  \endfirsthead
  \toprule
  Dimension & Extracted information \\
  \midrule
  \endhead
  Bibliographic identity & Title, authors, year, venue, DOI or stable URL,
    publication type, study family, evidence tier. \\
  Objective and task & Research objective; family, library, class, version,
    SBOM, or vulnerability granularity. \\
  Inputs and assumptions & Required APK artifacts, reference corpus,
    metadata, training data, and coverage assumptions. \\
  Segmentation & Package boundaries, class hierarchy, dependency graph,
    clustering, or other module-decoupling method. \\
  Evidence & Names, APIs, opcodes, control flow, call/dependency graphs,
    literals, resources, native code, metadata, learned representations, and
    cross-component relations. \\
  Decision method & Hashing, similarity, graph matching, ranking,
    classification, probability, learned model, or evidence fusion. \\
  Uncertainty behavior & Threshold, confidence, competing versions,
    abstention, partial matching, provenance, and conflict handling. \\
  Transformations & Renaming, package flattening, shrinking, optimization,
    control-flow rewriting, method inlining, repackaging, multidex, partial
    inclusion, native code, and dynamic loading. \\
  Evaluation & Dataset, ground truth, evaluation unit, metrics, baselines,
    obfuscation setting, and independent benchmark results. \\
  Operational support & Artifact availability, runtime, scalability, SBOM
    format, vulnerability mapping, and transitive-dependency handling. \\
  Limitations and relevance & Threats to validity, negative findings, and
    adoption, adaptation, or rejection relevance to the framework. \\
  \bottomrule
\end{longtable}

Unreported information is recorded as \texttt{NR}; it is not inferred from
neighboring fields or from another publication in the study family.

\section{Quality Assessment}

Tier A publications receive a score of 1, 0.5, or 0 for each question:

\begin{enumerate}
  \item Is the research objective clear?
  \item Is the method described well enough to understand or reproduce?
  \item Are the dataset and ground truth described?
  \item Are the task unit and metrics explicit?
  \item Are transformation or obfuscation conditions stated?
  \item Are identity and exact-version claims separated?
  \item Is an implementation, dataset, or artifact available?
  \item Are limitations or threats to validity acknowledged?
\end{enumerate}

The score controls synthesis weight rather than automatic inclusion.
Historically important studies may remain in the map with an explicit
reporting-quality limitation.

\subsection{Completed Quality Assessment}

\Cref{tab:review-quality-scores} records the completed assessment for the
15 current Tier A studies. The vector lists Q1--Q8 in the order defined above;
totals are out of eight. Scores reflect reporting and reproducibility in the
inspected publication, not the scientific importance of the work.

\begingroup
\scriptsize
\setlength{\LTpre}{0.5em}
\setlength{\LTpost}{0.5em}
\begin{longtable}{p{0.07\textwidth}p{0.33\textwidth}p{0.37\textwidth}p{0.10\textwidth}}
  \caption{Reporting-quality assessment for current Tier A studies}
  \label{tab:review-quality-scores}\\
  \toprule
  ID & Study & Q1--Q8 score vector & Total \\
  \midrule
  \endfirsthead
  \toprule
  ID & Study & Q1--Q8 score vector & Total \\
  \midrule
  \endhead
  S01 & LibScout & 1, 1, 1, 1, 1, 1, 1, 1 & 8.0 \\
  S02 & LibD & 1, 1, 0.5, 1, 0.5, 0.5, 0.5, 1 & 6.0 \\
  S03 & LibPecker & 1, 1, 1, 1, 1, 0.5, 0.5, 1 & 7.0 \\
  S04 & ORLIS & 1, 1, 1, 1, 1, 0.5, 0.5, 1 & 7.0 \\
  S05 & LibID & 1, 1, 1, 1, 1, 1, 0.5, 1 & 7.5 \\
  S06 & ASE 2020 assessment & 1, 1, 1, 1, 1, 1, 1, 1 & 8.0 \\
  S07 & Systematic assessment & 1, 1, 1, 1, 1, 1, 0.5, 1 & 7.5 \\
  S08 & ATVHUNTER & 1, 1, 1, 1, 1, 1, 0.5, 1 & 7.5 \\
  S09 & LibScan & 1, 1, 1, 1, 1, 1, 0.5, 1 & 7.5 \\
  S10 & Ramirez thesis & 1, 1, 1, 1, 0.5, 1, 0.5, 1 & 7.0 \\
  S11 & LibHunter & 1, 1, 1, 1, 1, 1, 0.5, 1 & 7.5 \\
  S12 & LibAttention & 1, 1, 1, 1, 1, 1, 0.5, 1 & 7.5 \\
  S13 & SAD & 1, 1, 1, 1, 1, 1, 0.5, 1 & 7.5 \\
  S17 & TPL-Benchmark & 1, 1, 1, 1, 1, 1, 1, 1 & 8.0 \\
  S18 & UniBOM & 1, 1, 0.5, 0.5, 0.5, 0.5, 1, 0.5 & 5.5 \\
  \bottomrule
\end{longtable}
\endgroup

\section{Synthesis Rules}

\begin{itemize}
  \item Narrative and taxonomy-based synthesis is the default.
  \item Incompatible datasets and metrics are not pooled into a common
    leaderboard.
  \item Library-level, class-level, and version-level results remain separate.
  \item Original-paper results remain separate from independent benchmark
    results.
  \item Obfuscation and optimization conditions must accompany performance
    claims.
  \item Artifact availability is assessed independently of reported accuracy.
  \item Tier B records cannot independently support precise numerical claims.
  \item Identity authority, version authority, and vulnerability authority are
    separate.
  \item Explicit ambiguity or abstention is preferred to unsupported exactness.
\end{itemize}

\section{Selected Studies}

\begin{longtable}{p{0.07\textwidth}p{0.55\textwidth}p{0.10\textwidth}p{0.18\textwidth}}
  \caption{Frozen seed corpus and current evidence status}
  \label{tab:review-seed-corpus}\\
  \toprule
  ID & Publication or study label & Year & Status \\
  \midrule
  \endfirsthead
  \toprule
  ID & Publication or study label & Year & Status \\
  \midrule
  \endhead
  S01 & LibScout \cite{backes2016libscout} & 2016 & Tier A \\
  S02 & LibD \cite{li2017libd} & 2017 & Tier A \\
  S03 & LibPecker \cite{zhang2018libpecker} & 2018 & Tier A \\
  S04 & ORLIS \cite{wang2018orlis} & 2018 & Tier A \\
  S05 & LibID \cite{zhang2019libid} & 2019 & Tier A \\
  S06 & Automated Third-Party Library Detection: Are We There Yet?
    \cite{zhan2020assessment} & 2020 & Tier A \\
  S07 & Systematic Assessment of Android TPL Detection Tools
    \cite{zhan2021systematic} & 2021 & Tier A \\
  S08 & ATVHUNTER \cite{zhan2021atvhunter} & 2021 & Tier A \\
  S09 & LibScan \cite{wu2023libscan} & 2023 & Tier A \\
  S10 & Empirical Study of TPL Detection in iOS and Android
    \cite{ramirez2024empirical} & 2024 & Tier A \\
  S11 & LibHunter / Code-Optimization Impact Study
    \cite{xie2024libhunter} & 2024 & Tier A \\
  S12 & LibAttention \cite{pan2025libattention} & 2025 & Tier A \\
  S13 & SAD \cite{zhou2025sad} & 2025 & Tier A \\
  S14 & LibRadar & 2016 & Tier B \\
  S15 & Libra & 2023 & Tier B \\
  S16 & LibMD & 2025 & Tier B \\
  S17 & TPL-Benchmark / Revisiting TPL Detection
    \cite{gu2025tplbenchmark} & 2025 & Tier A \\
  S18 & UniBOM \cite{safronov2025unibom} & 2025 & Tier A \\
  S19 & LibScope / Towards Global Matches & 2026 & Tier B \\
  S20 & ObfDetector / Advanced-Obfuscation Robustness Study & 2026 & Tier B \\
  \bottomrule
\end{longtable}

The seed corpus contains 20 publication units: 15 Tier A full texts and five
Tier B supplementary records. This count is a protocol checkpoint, not a claim
that all 20 have equal evidential authority. A Tier B record is promoted only
after full-text verification, bibliographic checking, structured extraction,
and quality assessment.

\section{Audit and Change Control}

The current extraction records are maintained under
\texttt{papers/data/}. Each search refresh must append the date, sources,
query variants, raw and deduplicated result counts, inclusion and exclusion
decisions, tier changes, and taxonomy effects to the master-thesis session log.
Eligibility, evidence tiers, and synthesis rules may be amended only through a
dated protocol change note; silent changes are prohibited.

\section{Protocol Change Notes}

\subsection{Version 1.0a --- Screening Clarification}

The full-text screening stage is complete only for records with a locally
verified full text. The frozen seed corpus predates the protocol, so its
original database-specific query counts cannot be reconstructed as an
exhaustive selection flow.

\subsection{Version 1.0b --- Full-Text Promotions}

On 21 June 2026, TPL-Benchmark and UniBOM were inspected locally, checked
bibliographically, extracted under the protocol schema, and quality-scored.
They were promoted from Tier B to Tier A. TPL-Benchmark supplies direct
Android benchmark evidence; UniBOM supplies adjacent binary/SBOM workflow
evidence and is not used to support Android Maven exact-version claims.
